\section{Vigen\`{e}re Cipher}
\label{sec:vigenere}

\subsection{Cipher Abjad-Majemuk}

Vigen\`{e}re cipher termasuk kelompok \crypt{cipher abjad-majemuk} (\textit{polyalphabetic cipher}). Perbedaannya dengan cipher abjad-tunggal \cite{munir_slides_itb}:
\begin{itemize}
  \item \textbf{cipher abjad-tunggal}: satu kunci untuk semua huruf plainteks (contoh: Caesar cipher);
  \item \textbf{cipher abjad-majemuk}: setiap huruf dapat dienkripsi dengan kunci yang berbeda (contoh: Vigen\`{e}re cipher).
\end{itemize}
Cipher abjad-majemuk dibuat untuk mengatasi kelemahan cipher abjad-tunggal yang mudah dipecahkan dengan analisis frekuensi.

Bentuk umum cipher abjad-majemuk adalah sebagai berikut. Kunci $K = k_1 k_2 \ldots k_m$ dengan $m$ panjang kunci. Karena kunci biasanya lebih pendek daripada plainteks $P = p_1 p_2 \ldots p_m p_{m+1} \ldots p_{2m} \ldots$, kunci diulang secara periodik:
\[
C = E_K(P) = f_{k_1}(p_1)\, f_{k_2}(p_2) \ldots f_{k_m}(p_m)\, f_{k_1}(p_{m+1})\, f_{k_2}(p_{m+2}) \ldots
\]
dengan $f(\cdot)$ fungsi enkripsi cipher abjad-tunggal. Untuk $m = 1$, cipher-nya ekivalen dengan cipher abjad-tunggal. Misalnya dengan kunci \texttt{sony} ($m = 4$), plainteks dipasangkan dengan kunci yang diulang:
\begin{center}
\ttfamily
\begin{tabular}{ll}
  \rmfamily Plainteks: & thisplaintexthissecret \\
  \rmfamily Kunci: & sonysonysonysonysonyso
\end{tabular}
\end{center}
Setiap huruf plainteks dienkripsi dengan huruf kunci di bawahnya, misalnya dengan Caesar cipher: $(t + s) \bmod 26$, $(h + o) \bmod 26$, $(i + n) \bmod 26$, dan seterusnya.

\subsection{Sejarah Singkat}

\begin{figure}[htbp]
  \centering
  \includegraphics[width=0.22\textwidth]{bab-04/potret-vigenere}
  \caption{Blaise de Vigen\`{e}re (1523--1596), diplomat dan kriptografer Prancis \cite{munir_slides_itb}.}
  \label{fig:potret-vigenere}
\end{figure}

Penemu cipher ini sebenarnya Giovan Battista Bellaso, yang menggambarkannya pada tahun 1553 dalam bukunya \textit{La Cifra del Sig. Giovan Battista Bellaso}. Cipher ini kemudian disempurnakan dan dipublikasikan oleh diplomat sekaligus kriptolog Prancis, Blaise de Vigen\`{e}re (Gambar~\ref{fig:potret-vigenere}), pada tahun 1586. Pada abad ke-19 banyak orang mengira Vigen\`{e}re adalah penemunya, sehingga cipher ini dikenal luas sebagai Vigen\`{e}re cipher \cite{munir_slides_itb}.

Vigen\`{e}re cipher dipakai oleh Tentara Konfederasi pada Perang Saudara Amerika (1861--1865), padahal saat itu cipher ini sudah dipecahkan secara diam-diam oleh Charles Babbage. Friedrich Kasiski kemudian mempublikasikan metode pemecahannya pada tahun 1863. Kasiski menguraikan langkah-langkah untuk menemukan \textbf{panjang kunci}, bukan langsung huruf-huruf kuncinya.

\subsection{Bujursangkar Vigen\`{e}re}

Enkripsi dan dekripsi dapat dilakukan dengan tabel bujursangkar Vigen\`{e}re (\textit{Vigen\`{e}re square}, Gambar~\ref{fig:bujursangkar-vigenere}). Setiap baris ke-$i$ (untuk huruf kunci ke-$i$) berisi huruf-huruf cipherteks hasil Caesar cipher dengan kunci $k = i$; artinya baris $i$ adalah alfabet yang digeser sejauh $i$ ke kanan. Huruf cipherteks diperoleh dari titik potong kolom huruf plainteks dan baris huruf kunci.

\begin{figure}[htbp]
  \centering
  \includegraphics[width=0.85\textwidth]{bab-04/bujursangkar-vigenere}
  \caption{Bujursangkar Vigen\`{e}re: enkripsi huruf plainteks \texttt{T} dengan huruf kunci \texttt{s} menghasilkan \texttt{L} \cite{munir_slides_itb}.}
  \label{fig:bujursangkar-vigenere}
\end{figure}

Tanpa bujursangkar pun, enkripsi dapat dihitung secara aritmetika: setiap huruf dienkripsi dengan Caesar cipher yang kuncinya berubah-ubah mengikuti huruf kunci. Untuk huruf plainteks $p_j$ dan huruf kunci $k_i$ yang bersesuaian:
\[
\text{Enkripsi: } c_j = (p_j + k_i) \bmod 26 \qquad
\text{Dekripsi: } p_j = (c_j - k_i) \bmod 26
\]

\begin{contoh}
Kunci \texttt{soreang} (dalam angka: 18, 14, 17, 4, 0, 13, 6). Kunci diulang sepanjang plainteks:
\begin{center}
\ttfamily
\begin{tabular}{ll}
  \rmfamily Plainteks: & terimapesanangulaikambing \\
  \rmfamily Kunci: & soreangsoreangsoreangsore \\
  \rmfamily Cipherteks: & LSIMMNVWGRRAAMMZRMKNSTWEK
\end{tabular}
\end{center}
Contoh perhitungan: $(\texttt{t} + \texttt{s}) \bmod 26 = (19 + 18) \bmod 26 = 11 = \texttt{L}$ dan $(\texttt{e} + \texttt{o}) \bmod 26 = (4 + 14) \bmod 26 = 18 = \texttt{S}$, dan seterusnya.
\end{contoh}

\subsection{Kelebihan: Histogram yang Datar}

Pada Vigen\`{e}re cipher, huruf plainteks yang sama tidak selalu dienkripsi menjadi huruf cipherteks yang sama; hasilnya bergantung pada huruf kunci. Pada contoh di atas, huruf plainteks \texttt{a} dienkripsi menjadi \texttt{A}, \texttt{N}, atau \texttt{R}, sedangkan huruf cipherteks \texttt{M} dapat merepresentasikan huruf plainteks \texttt{g}, \texttt{i}, \texttt{m}, atau \texttt{u}. Inilah karakteristik cipher abjad-majemuk: setiap huruf cipherteks memiliki banyak kemungkinan huruf plainteks.

Akibatnya, distribusi huruf di dalam cipherteks mendekati seragam (\textit{uniform}) dan histogramnya cenderung datar (\textit{flat}), seperti Gambar~\ref{fig:histogram-flat}. Histogram yang datar menyulitkan kriptanalis memecahkan cipherteks dengan analisis frekuensi biasa.

\begin{figure}[htbp]
  \centering
  \includegraphics[width=0.8\textwidth]{bab-04/histogram-flat}
  \caption{Histogram frekuensi huruf yang datar (kiri, khas cipher abjad-majemuk) dan yang tidak datar (kanan, khas bahasa alami dan cipher abjad-tunggal) \cite{munir_slides_itb}.}
  \label{fig:histogram-flat}
\end{figure}

\subsection{Extended Vigen\`{e}re Cipher}

Vigen\`{e}re cipher untuk 26 huruf dapat diperluas untuk mengenkripsi 256 karakter ASCII, sehingga dapat dipakai untuk sembarang berkas (teks maupun biner):
\[
c_j = (p_j + k_i) \bmod 256 \qquad p_j = (c_j - k_i) \bmod 256
\]

\begin{catatan}
\textbf{Kakas daring.} Vigen\`{e}re cipher dapat dicoba di \url{https://cryptii.com/pipes/vigenere-cipher}. Kakas \url{https://www.boxentriq.com/code-breaking/vigenere-cipher} bahkan menyediakan fitur \textit{auto-solve} yang memecahkan cipherteks secara otomatis.
\end{catatan}

\subsection{Implementasi dalam Python dan C}

Kode~\ref{lst:vigenere-py} dari slide kuliah memperlihatkan bahwa Vigen\`{e}re hanyalah Caesar cipher dengan kunci $k$ yang diambil bergiliran dari huruf-huruf kunci. Indeks kunci hanya maju ketika karakter yang diproses berupa huruf.

\begin{pycode}[caption={Vigen\`{e}re cipher dalam Python \cite{munir_slides_itb}}, label={lst:vigenere-py}]
def Vigenere_cipher_encrypt(plainteks, kunci):
    cipherteks = ""
    kunci = kunci.lower()
    indeks_kunci = 0
    for char in plainteks:
        if char.isalpha():        # Hanya memproses huruf alfabet saja
            start = ord('A') if char.isupper() else ord('a')
            k = ord(kunci[indeks_kunci]) - ord('a')
            c = (ord(char) - start + k) % 26  # Enkripsi dengan Caesar cipher
            c = c + start         # Kembalikan ke posisi semula
            cipherteks = cipherteks + chr(c)
            indeks_kunci = (indeks_kunci + 1) % len(kunci)
        else:
            cipherteks = cipherteks + char  # Bukan huruf: tidak dienkripsi
    return cipherteks

def Vigenere_cipher_decrypt(cipherteks, kunci):
    plainteks = ""
    kunci = kunci.lower()
    indeks_kunci = 0
    for char in cipherteks:
        if char.isalpha():
            start = ord('A') if char.isupper() else ord('a')
            k = ord(kunci[indeks_kunci]) - ord('a')
            c = (ord(char) - start - k) % 26  # Dekripsi dengan Caesar cipher
            c = c + start
            plainteks = plainteks + chr(c)
            indeks_kunci = (indeks_kunci + 1) % len(kunci)
        else:
            plainteks = plainteks + char
    return plainteks

print(Vigenere_cipher_encrypt("terimapesanangulaikambing", "soreang"))
# lsimmnvwgrraammzrmknstwek
\end{pycode}

\begin{ccode}[caption={Vigen\`{e}re cipher dalam C}, label={lst:vigenere-c}]
#include <stdio.h>
#include <string.h>
#include <ctype.h>

void vigenere_encrypt(const char *pt, const char *key, char *ct) {
    int klen = strlen(key), j = 0, i;
    for (i = 0; pt[i] != '\0'; i++) {
        if (isalpha((unsigned char)pt[i])) {
            char base = isupper((unsigned char)pt[i]) ? 'A' : 'a';
            int k = toupper((unsigned char)key[j % klen]) - 'A';
            ct[i] = base + (pt[i] - base + k) % 26;
            j++;                      /* kunci maju hanya untuk huruf */
        } else {
            ct[i] = pt[i];
        }
    }
    ct[i] = '\0';
}

void vigenere_decrypt(const char *ct, const char *key, char *pt) {
    int klen = strlen(key), j = 0, i;
    for (i = 0; ct[i] != '\0'; i++) {
        if (isalpha((unsigned char)ct[i])) {
            char base = isupper((unsigned char)ct[i]) ? 'A' : 'a';
            int k = toupper((unsigned char)key[j % klen]) - 'A';
            pt[i] = base + (ct[i] - base - k + 26) % 26;
            j++;
        } else {
            pt[i] = ct[i];
        }
    }
    pt[i] = '\0';
}

int main(void) {
    char enc[64], dec[64];
    vigenere_encrypt("terimapesanangulaikambing", "soreang", enc);
    printf("%s\n", enc);   /* lsimmnvwgrraammzrmknstwek */
    vigenere_decrypt(enc, "soreang", dec);
    printf("%s\n", dec);   /* terimapesanangulaikambing */
    return 0;
}
\end{ccode}

\subsection{Pengayaan: Kriptanalisis dengan Metode Kasiski}

\begin{figure}[htbp]
  \centering
  \includegraphics[width=0.2\textwidth]{bab-04/potret-kasiski}
  \caption{Friedrich Kasiski (1805--1881), perwira Prusia yang mempublikasikan pemecahan Vigen\`{e}re cipher pada tahun 1863 \cite{munir_slides_itb}.}
  \label{fig:potret-kasiski}
\end{figure}

Metode Kasiski (Gambar~\ref{fig:potret-kasiski}) tidak langsung menemukan kunci, tetapi membantu memperkirakan \textbf{panjang} kunci. Metode ini memanfaatkan fakta bahwa bahasa Inggris tidak hanya mengandung perulangan huruf, tetapi juga perulangan pasangan dan tripel huruf seperti TH, THE, dan EN. Perulangan kelompok huruf ini dapat menghasilkan kriptogram yang berulang pula.

\begin{contoh}
Kata \texttt{crypto} muncul dua kali dengan jarak 16 huruf.
\begin{center}
\ttfamily\small
\begin{tabular}{ll}
  \rmfamily Plainteks: & cryptoisshortforcryptography \\
  \rmfamily Kunci: & abcdabcdabcdabcdabcdabcdabcd \\
  \rmfamily Cipherteks: & \textbf{CSASTP}KVSIQUTGQU\textbf{CSASTP}IUAQJB \\[4pt]
  \rmfamily Kunci: & abcdefabcdefabcdefabcdefabcd \\
  \rmfamily Cipherteks: & CSASXTITUKSWTGQUGWYQVRKWAQJB
\end{tabular}
\end{center}
Dengan kunci \texttt{abcd} (panjang 4), kedua \texttt{crypto} dienkripsi menjadi kriptogram yang sama, \texttt{CSASTP}, karena jaraknya (16) merupakan kelipatan 4. Dengan kunci \texttt{abcdef} (panjang 6), kriptogramnya berbeda karena 16 bukan kelipatan 6.
\end{contoh}

Secara intuitif: jika jarak antara dua \textit{string} yang berulang di dalam plainteks merupakan kelipatan panjang kunci, \textit{string} tersebut akan muncul sebagai kriptogram yang sama di dalam cipherteks. Tujuan metode Kasiski adalah mencari dua atau lebih kriptogram berulang untuk menentukan panjang kunci, dengan langkah-langkah:
\begin{enumerate}
  \item Temukan semua kriptogram yang berulang di dalam cipherteks (pesan yang panjang biasanya mengandung kriptogram berulang).
  \item Hitung jarak antara kriptogram-kriptogram yang berulang.
  \item Hitung semua faktor pembagi setiap jarak; faktor pembagi menyatakan panjang kunci yang mungkin.
  \item Tentukan irisan himpunan faktor pembagi tersebut. Nilai yang muncul pada semua himpunan kemungkinan besar adalah panjang kunci.
\end{enumerate}

\begin{contoh}
Cipherteks berikut memuat dua kriptogram berulang:
\begin{center}
\texttt{\textbf{DYDUXRMH}TVDV\underline{NQD}QNW\textbf{DYDUXRMH}ARTJGW\underline{NQD}}
\end{center} Jarak antara dua \texttt{DYDUXRMH} adalah 18, dengan faktor pembagi $\{2, 3, 6, 9, 18\}$. Jarak antara dua \texttt{NQD} adalah 20, dengan faktor pembagi $\{2, 4, 5, 10, 20\}$. Irisan kedua himpunan adalah $\{2\}$, sehingga panjang kunci kemungkinan besar 2.
\end{contoh}

Setelah panjang kunci $p$ diketahui, huruf-huruf kunci dapat dicari dengan \textit{exhaustive key search} (paling banyak $26^p$ percobaan), tetapi analisis frekuensi jauh lebih sangkil:
\begin{enumerate}
  \item Setiap huruf yang berjarak $p$ pasti dienkripsi dengan huruf kunci yang sama. Kelompokkan huruf-huruf berjarak $p$ sehingga diperoleh $p$ ``pesan'', masing-masing dienkripsi dengan Caesar cipher (cipher abjad-tunggal).
  \item Pecahkan setiap ``pesan'' dengan analisis frekuensi.
  \item Susun huruf-huruf kunci dari hasil langkah 2, atau terka kata yang membantu memecahkan cipherteks.
\end{enumerate}

\begin{figure}[htbp]
  \centering
  \includegraphics[width=\textwidth]{bab-04/kasiski-ljv}
  \caption{Cipherteks Vigen\`{e}re dengan enam kemunculan kriptogram \texttt{LJV} \cite{munir_slides_itb}.}
  \label{fig:kasiski-ljv}
\end{figure}

\begin{contoh}
Pada cipherteks Gambar~\ref{fig:kasiski-ljv}, kriptogram \texttt{LJV} muncul enam kali. Jarak antarkemunculannya berturut-turut 15, 15, 15, 10, dan 10. Faktor pembagi 15 adalah $\{3, 5, 15\}$ dan faktor pembagi 10 adalah $\{2, 5, 10\}$; irisannya $\{5\}$, sehingga panjang kunci diperkirakan 5.

Cipherteks lalu dikelompokkan per jarak 5, dimulai dari huruf pertama, kedua, dan seterusnya:
\begin{center}
\small
\begin{tabular}{cl@{\hspace{2em}}c}
  \toprule
  Kelompok & Pesan & Huruf tersering \\
  \midrule
  1 & \texttt{LSLLMFLYJLVLLLYKWV} & L \\
  2 & \texttt{JTQJPAJYQJTJKJJREH} & J \\
  3 & \texttt{VNMVKVVVLVRVEVVFFZ} & V \\
  4 & \texttt{BEEMAADNNHNCTHSTUN} & N \\
  5 & \texttt{QZDAUTAFPKFMAUFTXP} & A \\
  \bottomrule
\end{tabular}
\end{center}

Triplet tersering dalam bahasa Inggris adalah THE, sehingga \texttt{LJV} diterka sebagai \texttt{THE}. Huruf kunci kelompok lain diterka dan diuji coba dengan cara serupa. Ingat bahwa nilai numerik huruf kunci adalah besar pergeseran Caesar cipher:
\begin{center}
\small
\begin{tabular}{cccc}
  \toprule
  Kelompok & Huruf plainteks & Huruf cipherteks & Huruf kunci \\
  \midrule
  1 & T & L & S (= 18) \\
  2 & H & J & C (= 2) \\
  3 & E & V & R (= 17) \\
  4 & N & N & A (= 0) \\
  5 & O & A & M (= 12) \\
  \bottomrule
\end{tabular}
\end{center}
Jadi kuncinya \texttt{SCRAM}, dan cipherteks terdekripsi menjadi \texttt{THEBE ARWEN TOVER THEMO UNTAI NYEAH \ldots}. Setelah \textit{post-processing}: \textit{THE BEAR WENT OVER THE MOUNTAIN YEAH, THE DOG WENT ROUND THE HYDRANT, THE CAT INTO THE HIGHEST SPOT HE COULD FIND}.
\end{contoh}

\begin{catatan}
Panjang kunci juga dapat diperkirakan dengan \crypt{index of coincidence} yang dikembangkan William Friedman (1920-an): jika cipherteks dibagi menjadi $p$ kelompok dan $p$ adalah panjang kunci yang benar, indeks koinsidensi setiap kelompok mendekati nilai bahasa asli (sekitar 0,065 untuk bahasa Inggris); jika salah, nilainya mendekati teks acak (sekitar 0,038).
\end{catatan}

\subsection{Pengayaan: Varian Vigen\`{e}re Cipher}

Untuk menahan serangan metode Kasiski, dibuat beberapa varian Vigen\`{e}re cipher \cite{munir_slides_itb}:
\begin{enumerate}
  \item \textbf{Full Vigen\`{e}re cipher}. Setiap baris bujursangkar bukan pergeseran alfabet, melainkan permutasi acak huruf-huruf alfabet. Tabel ini harus dirahasiakan; proses enkripsi dan dekripsi tetap sama seperti Vigen\`{e}re standar.
  \item \textbf{Auto-key Vigen\`{e}re cipher}. Jika kunci lebih pendek daripada plainteks, kunci disambung dengan plainteks itu sendiri, sehingga tidak ada pengulangan periodik. Contoh dengan pesan \texttt{negara penghasil minyak mentah di dunia} dan kunci \texttt{INDO}:
  \begin{center}
  \ttfamily\small
  \begin{tabular}{ll}
    \rmfamily Plainteks: & negarapenghasilminyakmentahdidunia \\
    \rmfamily Kunci: & INDONEGARAPENGHASILMINYAKMENTAHDID \\
    \rmfamily Cipherteks: & VRJOEEVEEGWEFOSMAVJMSZCNDMLQBDBQQD
  \end{tabular}
  \end{center}
  \item \textbf{Running-key Vigen\`{e}re cipher}. Kunci adalah \textit{string} sangat panjang yang diambil dari teks bermakna, misalnya naskah proklamasi, Pembukaan UUD 1945, atau terjemahan kitab suci. Untuk pesan yang sama:
  \begin{center}
  \ttfamily\small
  \begin{tabular}{ll}
    \rmfamily Plainteks: & negarapenghasilminyakmentahdidunia \\
    \rmfamily Kunci: & KERAKYATANYANGDIPIMPINOLEHHIKMATPE \\
    \rmfamily Cipherteks: & XIXABYPXNTFAFOOUXVKPSZSYXHOLSPUGXE
  \end{tabular}
  \end{center}
\end{enumerate}
